%% kapa: exam=1 round=22 year=2011 period=1 form=A subject=civil_law subject_ko=민법 seq=19 number=19
%% source: Q-Net 감정평가사 기출문제 (cst00309) artlSeq=5036790 file=extracted/1008/5036790/2011년 제22회 감정평가사 1차시험 기출문제/1교시.hwp
%% conversion: kordoc 4.24.51 -> tools/gen1.py (auto); review=proofread-sonnet; flags=none
\begin{question}{19}
甲은 乙에 대한 채권을 담보하기 위해 乙소유의 X부동산, 丙소유의 Y부동산에 각각 1번 저당권을 취득하였다. 甲은 Y부동산에 대해 담보권을 실행하여 그 경매대금으로부터 채권 전액을 배당받았다. 다음 설명으로 옳지 않은 것은? (다툼이 있으면 판례에 의함)
\begin{choices}
  \item 丙은 X부동산에 대한 1번 저당권을 대위취득한다.
  \item 丙은 乙에 대하여 구상권을 취득한다.
  \item Y부동산의 2번 저당권자는 X부동산의 1번 저당권에 대하여 물상대위를 할 수 있다.
  \item 乙은 甲에게 피담보채무가 소멸되었음을 이유로 X부동산의 1번 저당권의 말소등기를 청구할 수 있다.
  \item Y부동산의 2번 저당권자는 자신의 채무자 丙을 대위하여 甲에게 X부동산의 1번 저당권 이전의 부기등기를 할 것을 청구할 수 있다.
\end{choices}
\end{question}
\answer{4}
